
\subsection{Spatial Relation Extraction}
\label{sec:spatial_relation_extraction}

% Task definition
%  - Motivation & Importance
%  - Formal Task Definition
%  - Subtasks, if have:

% Part 2: Method taxonomy.
% modality
%

% Task Definition
%  - Motivation & Importance
% ,liu2024edge，liu2025boost
Spatial Relation Extraction (SRE) identifies the relative spatial configuration of objects, such as ``table-left-of-chair'' or ``car-behind-bus''.
It turns detections or segmentations into spatial relation knowledge that feeds reasoning, scene graph generation~\cite{xu2025llava,li2024pixels,wu2025universal}, autonomous driving~\cite{malawade2022roadscene2vec,zhou2023pit,liu2024edge,liu2025boost}, and embodied AI~\cite{SongBTTSB25,zhang2025vlabench}.
It is narrower than Scene Graph Generation (SGG), which runs the full pipeline from detection to relation prediction and includes non-spatial predicates (``person-holding-book''); here SRE covers only geometric, topological, and projective spatial predicates.


%  - Formal Task Definition
\subsubsection{Task Definition}
Given a visual scene (a 2D image, a video, a point cloud, or an RGB-D stream) and two target object instances specified by tags, boxes, or masks, SRE outputs their spatial relationship either as a triplet $(s, p, o)$ with $p$ drawn from a spatial predicate vocabulary, or as a natural-language sentence describing the relation and its context.
Three subtasks differ in output granularity. \textbf{Visual Spatial Classification (VSC)}~\cite{yang2019spatialsense, zhao2024acm} predicts a predicate from a closed label set (``on'', ``in'', ``next to'') or verifies whether a given predicate holds for the object pair. \textbf{Visual Spatial Description (VSD)}~\cite{zhao2022visual, zhao2023generating} generates a free-form sentence that can carry attributes and context beyond the bare predicate. \textbf{Quantitative spatial understanding}~\cite{cheng2024spatialrgpt,chen2024spatialvlm} replaces qualitative predicates with metric quantities (distances, absolute sizes, 3D orientations), which navigation and manipulation require. The field has moved through these three in order, from closed-set classification to generation to metric quantification.


% Part 2: Method taxonomy.
% (Focus on giving a comprehensive but structured overview; no need to cover every paper in depth.) 规则+学习
\subsubsection{Methodology}

Figure~\ref{fig:p3_s6_3_1_sre_overview} summarizes the task and the three method families surveyed below.

\begin{figure}[!t]
    \centering
    \includegraphics[width=\linewidth]{figures/part3/section6/spatial_relation_extraction.pdf}
    \caption{Spatial relation extraction. \textit{Top}: given a visual scene (image, video, point cloud, or RGB-D) and two target objects $(O_1, O_2)$, predict (a) a closed-set relation label (VSC), (b) a free-form description (VSD), or (c) a quantitative metric estimate. \textit{Bottom}: (i) rule-based, (ii) learning-based decoupled, and (iii) MLLM-based mechanisms with method timelines (chronological spacing is not proportional). Metric outputs require depth or scale information.}
    \label{fig:p3_s6_3_1_sre_overview}
\end{figure}

\paragraph{Rule-based approaches.}
Rule-based pipelines~\cite{malawade2022roadscene2vec,zhao2022visual,zhou2024hktsg} first run off-the-shelf perception (a 2D detector with a bird's-eye-view transform~\cite{liu2008bird}, or a 3D scene extractor~\cite{nie2020total3dunderstanding}) to obtain each object's 3D position, size, orientation, and appearance in a camera-centered frame whose vertical axis is aligned with gravity, and then apply deterministic geometric rules to the centroids to produce $(O_1, p, O_2)$ triplets~\cite{zhao2022visual}.
Let $(x_{O_i}, y_{O_i}, z_{O_i})$ denote the centroid of object $O_i$ in this shared camera-centered frame, with $x$ pointing right and $y$ pointing upward. Then
\begin{equation}
\begin{aligned}
\Delta x &= |x_{O_1} - x_{O_2}|, \quad
\Delta y = |y_{O_1} - y_{O_2}|, \quad
\Delta z = |z_{O_1} - z_{O_2}|, \\
&\Delta y > \tau \,\land\, y_{O_1} > y_{O_2} \;\Rightarrow\; \text{``above''}, \\
&\Delta x > \Delta z \,\land\, \Delta x > \tau \,\land\, x_{O_1} < x_{O_2} \;\Rightarrow\; \text{``left''}, \\
&\Delta x > \Delta z \,\land\, \Delta x > \tau \,\land\, x_{O_1} > x_{O_2} \;\Rightarrow\; \text{``right''},
\end{aligned}
\end{equation}
where $\tau$ is a small threshold (\eg, $0.2$\,m).
Such rules are cheap but brittle: 2D-to-3D back-projection misplaces occluded or oddly posed objects (a rear vehicle whose centroid lands ambiguously on the BEV plane), fixed rules cannot express functional or contextual relations (``person riding bicycle'', a book more likely on a shelf than under it), and every new domain needs its rules rewritten.


\paragraph{Learning-based approaches.}
Zhao \etal~\cite{zhao2022visual} decoupled spatial relation generation from generic captioning and trained it with a joint visual--spatial objective; 3DVSD~\cite{zhao2023generating} moved the same task into 3D scene modeling and added structure-aware sampling that holds up under occlusion.

\paragraph{MLLM-based approaches.}
Multimodal large language models carried the line from image-level predicates toward finer and more metric grounding.
SpatialVLM~\cite{chen2024spatialvlm} used monocular-depth-derived spatial QA supervision to train a vision-language model so the model answers metric queries such as ``How far is A from B?'', outperforming GPT-4V on spatial VQA.
SpatialBot~\cite{cai2024spatialbot} exposes pixel-level depth queries through a depth API and trains progressively, translating spatial reasoning into higher robotic grasping success.
SpatialRGPT~\cite{cheng2024spatialrgpt} sharpens granularity from the whole image to regions, accepting boxes or masks and auto-generating 8.7M region-level spatial samples through scene-subgraph sampling (S$^3$).
SpatialLLM~\cite{ma2025spatialllm} pairs a hybrid CLIP--DINOv2 encoder with staged data injection to hold qualitative and quantitative capabilities in one model.


% Part 3: Benchmarks & metrics.
\subsubsection{Datasets and Metrics}
Table~\ref{tab:sre_dataset} lists the benchmarks. They differ mainly in how relations were collected: manual labeling of robot-captured lab scenes in SpatialAwareRobot~\cite{wang2025spatial}; adversarial crowdsourcing in SpatialSense~\cite{yang2019spatialsense}, where annotators propose relations that fool a baseline and only those count, which makes the dataset deliberately hard; the merge of SpatialSense and Visual Genome~\cite{krishna2017visual} into VSD~\cite{zhao2022visual}, whose VSD-v1 version contains over 143K descriptions; the 34-preposition, English--French annotations of SpatialVOC2K~\cite{belz2018spatialvoc2k}; and the automatically generated, depth-aware SpaceSGG~\cite{xu2025llava}, released as 10K scene descriptions, 20K question--answer pairs, and 20K multi-turn conversations over 20K scenes.

\begin{table}[!t]
\centering
\caption{Benchmark datasets for spatial relation extraction. ``--'' marks an open vocabulary without a fixed category count.}
\label{tab:sre_dataset}
\fontsize{9}{11}\selectfont
\setlength{\tabcolsep}{3mm}
\begin{tabular}{lccccc}
\toprule
\textbf{Dataset} & \textbf{Scene Type} & \textbf{Modality} & \textbf{Annotation} & \textbf{Samples} & \textbf{Relation Categories} \\
\midrule
SpatialAwareRobot~\cite{wang2025spatial} & Robotics & Image & Manual & 1,000 & 7 \\
SpatialSense~\cite{yang2019spatialsense} & Generic & Image & Manual & 11,569 & 9 \\
VSD~\cite{zhao2022visual} & Generic & Image & Manual & 29,272  & -- \\
SpatialVOC2K~\cite{belz2018spatialvoc2k} & Generic & Image & Manual & 2,026  & 34 \\
SpaceSGG~\cite{xu2025llava} & Generic & Image & Automated & 20,000  & --  \\
\bottomrule
\end{tabular}
\end{table}

Each subtask has its own evaluation regime, and the metrics are not comparable across them.
VSC is scored as classification: SpatialSense~\cite{yang2019spatialsense} measures binary accuracy for given object--predicate pairs, with the driving-scene systems Roadscene2Vec~\cite{malawade2022roadscene2vec} and HKT-SG~\cite{zhou2024hktsg} reporting per-relation F1 or accuracy on their own data; results on SpatialSense's adversarial split lag well behind those on conventional crowdsourced data, which is the contextual-relation difficulty taken up under Challenges.
VSD is scored as generation: VSDNet~\cite{zhao2022visual} and 3DVSD~\cite{zhao2023generating} report BLEU-4, METEOR, and CIDEr on the VSD datasets, and LLaVA-SpaceSGG~\cite{xu2025llava} reports QA accuracy on SpaceSGG.
Quantitative understanding is the newest regime: SpatialVLM~\cite{chen2024spatialvlm} and SpatialRGPT~\cite{cheng2024spatialrgpt} evaluate metric estimates using human-annotated spatial questions and answers (SpatialVLM) or questions derived from ground-truth 3D cuboid annotations (SpatialRGPT-Bench), SpatialLLM~\cite{ma2025spatialllm} reports combined qualitative and quantitative VQA accuracy, and SpatialBot~\cite{cai2024spatialbot} adds robotic grasp success as an embodied proxy.
MLLM-based models dominate the qualitative leaderboards but still fall short on metric-scale reasoning, which motivates the unified representations discussed next.


\subsubsection{Challenges}
Spatial relations are not purely geometric. ``A cup on a table'' is plausible and ``a cup on a wall'' is not, yet the two differ only in object identity, so rule-based and purely geometric methods cannot capture functional or commonsense constraints; a model needs to know what objects are for and where they tend to be, not just their boxes or point coordinates.
Metric grounding is the second gap. MLLMs handle contextual and functional predicates well but judge distances, sizes, and angles poorly; injecting geometry and depth during pre- or post-training, or exposing dedicated spatial tools, is the current route, and the longer-term goal is a representation that holds discrete predicates (``near'') and continuous quantities (``2 meters away'') together, since navigation and manipulation need both.
Generalization is the third. SpatialSense~\cite{yang2019spatialsense} and VSD~\cite{zhao2022visual} cover a small slice of real scenes, and models trained on them degrade under new object appearances, scene complexity, lighting, and occlusion; SpatialVLM~\cite{chen2024spatialvlm} and its successors generalize better, but unseen and dynamic environments remain largely untested.

